\section{Giá trị lượng giác của một góc lượng giác}

\subsection{Đề 1}

\Opensolutionfile{ans}[ans/ans-0-B1]
\caulc
\begin{ex}%[1D1N2-1]%Câu 1%
	Trong các hệ thức sau đây, hệ thức nào đúng?
	\choice
	{\True $\sin^2x+\cos^2x=1$}
	{$\sin{x^2}+\cos{x^2}=1$}
	{$\sin 2x+\cos 2x=1$}
	{$\sin^2x+\cos{x^2}=1$}
	\loigiai{
	$\sin^2x+\cos^2x=1$.
}
	\end{ex}
	

	
	\begin{ex}%[1D1N2-2]
		Cho góc $\alpha $ thỏa mãn $ 2\pi <\alpha <\dfrac{5\pi}{2}$. Khẳng định nào sau đây {\bf sai}?
		\choice
		{\True $\tan\alpha <0$}
		{$\cot\alpha >0$}
		{$\sin\alpha >0$}
		{$\cos\alpha >0$}
		\loigiai
		{	Với $ 2\pi <\alpha <\dfrac{5\pi}{2}$ ta có $\sin\alpha >0$, $\cos\alpha >0$, $\tan\alpha >0$, $\cot\alpha >0$.
		}
		\end{ex}
		
	
		
		\begin{ex}%[1D1N2-2]
			Cho $ 0<\alpha <\dfrac{\pi}{2}$. Mệnh đề nào sau đây sai?
			\choice
			{$\sin (\alpha\,+\,\pi)<0$}
			{$\cos (\alpha\,+\,\pi)<0$}
			{\True $\tan (\pi\,-\,\alpha)>0$}
			{$\cot (\pi\,-\,\alpha)<0$}
			\loigiai{$\tan (\pi-\alpha)=-\tan\alpha <0\,\forall\,\alpha\in\left(0\,;\,\dfrac{\pi}{2}\right)$.}
		\end{ex}
		
	
			
			\begin{ex}%[1D1N2-2]
				Cho $ 0<\alpha <\dfrac{\pi}{2}$. Khẳng định nào sau đây đúng?
				\choice
				{\True $\tan\alpha >0,\cot\alpha >0$}
				{$\tan\alpha >0,\cot\alpha <0$}
				{$\tan\alpha <0,\cot\alpha <0$}
				{$\tan\alpha <0,\cot\alpha >0$}
				\loigiai
				{		Điểm biểu diễn cung lượng giác $\alpha$ nằm góc phần tư thứ nhất nên $\tan\alpha >0$, $\cot\alpha >0$.}
				\end{ex}
				
		
				
				\begin{ex}%[1D1H2-2]
					Tìm khẳng định {\bf sai} trong các khẳng định sau đây?
					\choice
					{$\tan 45^\circ <\tan 60^\circ $}
					{$\cos 45^\circ\le\sin 45^\circ $}
					{$\sin 60^\circ <\sin 80^\circ $}
					{\True $\cos 35^\circ >\cos 10^\circ $}
					\loigiai
					{Khi $\alpha\in\left(0^\circ\,;\,90^\circ\right)$ hàm $\cos\alpha $ là hàm giảm nên $\cos 35^\circ <\cos 10^\circ $ suy ra $\cos 35^\circ >\cos 10^\circ $ sai.}
				\end{ex}
				
			
				\begin{ex}%[1D1N2-1]
					Giá trị của $\tan\dfrac{\pi}{6}$ là
					\choice
					{\True $\dfrac{\sqrt{3}}{3}$}
					{$-\dfrac{\sqrt{3}}{3}$}
					{$\sqrt{3}$}
					{$-\sqrt{3}$}
					\loigiai{
					Trên nửa đường tròn đơn vị sao cho $\widehat{xOM}=\dfrac{\pi}{6}$.\\
					Khi đó $M\left(\dfrac{\sqrt{3}}{2}\,;\,\dfrac{1}{2}\right)$ nên $\tan\dfrac{\pi}{6}=\dfrac{\dfrac{1}{2}}{\dfrac{\sqrt{3}}{2}}=\dfrac{\sqrt{3}}{3}$.}
				\end{ex}
				
			
				\begin{ex}%[1D1H2-4]
					Hãy chọn kết quả sai trong các kết quả sau đây
					\choice
					{$\cos\left(-\alpha\right)=\cos\alpha $}
					{\True $\sin\left(\pi+\alpha\right)=\sin\alpha $}
					{$\tan\left(\pi-\alpha\right)=-\tan\alpha $}
					{$\cot\left(\dfrac{\pi}{2}-\alpha\right)=\tan\alpha $}
					\loigiai
					{Ta có $\sin\left(\pi+\alpha\right)=-\sin\alpha $ nên $\sin\left(\pi+\alpha\right)=\sin\alpha $ {\bf sai}.}
				\end{ex}
				
			
				
				\begin{ex}%[1D1H2-2]
					Giá trị của biểu thức $ S=3-\sin^290^\circ+2\cos^260^\circ-3\tan^245^\circ $ bằng
					\choice
					{$\dfrac{1}{2}$}
					{\True $-\dfrac{1}{2}$}
					{$ 1$}
					{$ 3$}
					\loigiai
					{Ta có $ S=3-\sin^290^\circ+2\cos^260^\circ-3\tan^245^\circ =3-1^2+2\cdot \left(\dfrac{1}{2}\right)^2-3\cdot 1^2=-\dfrac{1}{2}$.}
				\end{ex}
				
			
		
				\begin{ex}%[1D1H2-2]
					Cho góc $\alpha $ thỏa $\cot\alpha=\dfrac{3}{4}$ và $0^\circ<\alpha <90^\circ$. Khẳng định nào sau đây đúng?
					\choice
					{$\cos\alpha=\dfrac{4}{5}$}
					{\True $\sin\alpha=\dfrac{4}{5}$}
					{$\sin\alpha=-\dfrac{4}{5}$}
					{$\cos\alpha=-\dfrac{4}{5}$}
					\loigiai
					{Áp dụng công thức $ 1+\cot^2\alpha=\dfrac{1}{\sin^2\alpha}$ $\Leftrightarrow 1+\dfrac{9}{16}=\dfrac{1}{\sin^2\alpha}$ $\Leftrightarrow{\sin^2}\alpha=\dfrac{16}{25}$.\\
					Do $0^\circ<\alpha <90^\circ\Rightarrow\sin\alpha >0$ $\Rightarrow\sin\alpha=\dfrac{4}{5}$.}
				\end{ex}
				
			
							
							\begin{ex}%[1D1V2-2]
								Giá trị của $\tan\left(\alpha+\dfrac{\pi}{3}\right)$ bằng bao nhiêu khi $\sin\alpha=\dfrac{3}{5}$ $\left(\dfrac{\pi}{2}<\alpha <\pi\right)$?
								\choice
								{$\dfrac{-48+25\sqrt{3}}{11}$}
								{$\dfrac{48+25\sqrt{3}}{11}$}
								{$-\dfrac{48+25\sqrt{3}}{11}$}
								{\True $\dfrac{48-25\sqrt{3}}{11}$}
								\loigiai{
									Ta có $\sin^2\alpha+\cos^2\alpha=1$
									$\Rightarrow\cos\alpha=\pm\sqrt{1-\sin^2\alpha}$
									$\Rightarrow\cos\alpha=\pm\sqrt{1-\left(\dfrac{3}{5}\right)^2}=\pm\dfrac{4}{5}$.\\
									Mà $\dfrac{\pi}{2}<\alpha <\pi$ nên $\cos\alpha=-\dfrac{4}{5}$. Suy ra $\tan\alpha=\dfrac{-3}{4}$.\\
									Khi đó ta có $\tan\left(\alpha+\dfrac{\pi}{3}\right)=\dfrac{\tan\alpha+\tan\dfrac{\pi}{3}}{1-\tan\alpha \cdot \tan\dfrac{\pi}{3}}=\dfrac{\dfrac{-3}{4}+\sqrt{3}}{1-\left(\dfrac{-3}{4}\right)\cdot \sqrt{3}}=\dfrac{48-25\sqrt{3}}{11}$.}
							\end{ex}
							
		
								\begin{ex}%[1D1V2-2]
									Biết $\sin\alpha+\cos\alpha=\dfrac{\sqrt{3}}{2}$. Trong các kết quả sau, kết quả nào {\bf sai}?
									\choice
									{$\sin\alpha .\cos\alpha=-\dfrac{1}{8}$}
									{$\sin\alpha-\cos\alpha=\pm\dfrac{\sqrt{5}}{2}$}
									{\True $\sin^4\alpha+\cos^4\alpha=\dfrac{5}{4}$}
									{$\tan^2\alpha+\cot^2\alpha=62$}
									\loigiai
									{		$\begin{aligned}
											\sin\alpha+\cos\alpha=\dfrac{\sqrt{3}}{2}&\Rightarrow{\left(\sin\alpha+\cos\alpha\right)^2}=\left(\dfrac{\sqrt{3}}{2}\right)^2\\ 
											&\Leftrightarrow{\sin^2}\alpha+\cos^2\alpha+2\sin\alpha \cdot \cos\alpha=\dfrac{3}{4}\\ 
											&\Leftrightarrow 1+2\sin\alpha \cdot \cos\alpha=\dfrac{3}{4}\\ 
											&\Leftrightarrow\sin\alpha \cdot \cos\alpha=-\dfrac{1}{8}. 
										\end{aligned}$\\
										$\left(\sin\alpha-\cos\alpha\right)^2=\sin^2\alpha+\cos^2\alpha-2\sin\alpha \cdot \cos\alpha=1+\dfrac{1}{4}=\dfrac{5}{4}$ 
										$\Rightarrow\sin\alpha-\cos\alpha=\pm\dfrac{\sqrt{5}}{2}$.\\
										$\sin^4\alpha+\cos^4\alpha=\left(\sin^2\alpha+\cos^2\alpha\right)^2-2\sin^2\alpha \cdot \cos^2\alpha=1-2\cdot \left(-\dfrac{1}{8}\right)^2=\dfrac{31}{32}$.\\
										$\tan^2\alpha+\cot^2\alpha=\dfrac{\sin^2\alpha}{\cos^2\alpha}+\dfrac{\cos^2\alpha}{\sin^2\alpha}=\dfrac{\sin^4\alpha+\cos^4\alpha}{\sin^2\alpha \cdot \cos^2\alpha}=\dfrac{\dfrac{31}{32}}{\dfrac{1}{64}}=62$.
									}
								\end{ex}
								
	
								\begin{ex}%[1D1V2-2]
									Biểu thức $A=\sin^210^\circ+\sin^220^\circ+\cdots +\sin^2180^\circ $ có giá trị bằng
									\choice
									{$A=6$}
									{$A=8$}
									{$A=3$}
									{\True $A=9$}
									\loigiai
									{	Ta có $\sin\left(90^\circ+\alpha\right)=\cos\alpha $.\\
										Suy ra $\sin 100^\circ=\cos 10^\circ\Rightarrow{\sin^2}100^\circ=\cos^210^\circ $.\\
										Tương tự ta có $\sin^2110^\circ=\cos^220^\circ $, $\sin^2120^\circ=\cos^230^\circ$, $\sin^2130^\circ=\cos^240^\circ $,\\
										$\sin^2150^\circ=\cos^240^\circ $, $\sin^2160^\circ=\cos^270^\circ $, $\sin^2170^\circ=\cos^280^\circ $, $\sin^2180^\circ=\cos^290^\circ $.\\
										Vậy ta có $A=\left(\sin^210^\circ+\cos^210^\circ\right)+\left(\sin^220^\circ+\cos^220^\circ\right)+ \cdots +\left(\sin^290^\circ+\cos^290^\circ\right)$		\\								$\Rightarrow A=1+1+\cdots +1=9$.}
								\end{ex}
								
		

\Closesolutionfile{ans}
\indapan{6}{ans/ans-0-B1}
\cauds
\Opensolutionfile{ans}[ans/ans-0-B1-DS]



\begin{ex}%[1D1H2-2]
Cho $\alpha=\dfrac{\pi}{3}+k2\pi $ (biết $k\in\mathbb{Z}$). Các khẳng định  sau đây đúng hay sai? 
\choiceTF
{$\sin\alpha=-\dfrac{\sqrt{3}}{2}$}
{$\cos\alpha=-\dfrac{1}{2}$}
{\True $\tan\alpha=\sqrt{3}$}
{$\cot\alpha=-\dfrac{\sqrt{3}}{3}$}
\loigiai{
\begin{itemchoice}
\itemch  $ \sin\left(\dfrac{\pi}{3}+k 2\pi\right)=\sin\dfrac{\pi}{3}=\dfrac{\sqrt{3}}{2}$. 
\itemch  $ \cos\left(\dfrac{\pi}{3}+k 2\pi\right)=\cos\dfrac{\pi}{3}=\dfrac{1}{2}$. 
\itemch $ \tan\left(\dfrac{\pi}{3}+k 2\pi\right)=\tan\dfrac{\pi}{3}=\sqrt{3}$.
\itemch  $ \cot\left(\dfrac{\pi}{3}+k 2\pi\right)=\cot\dfrac{\pi}{3}=\dfrac{\sqrt{3}}{3}$
\end{itemchoice}
}
\end{ex}


\begin{ex}%[1D1H2-2]
Cho $0^{\circ}<\alpha<90^{\circ}$. Các khẳng định  sau đây đúng hay sai? 
\choiceTF
{\True $ A=\sin\left(\alpha+90^{^\circ}\right)>0$}
{\True $ B=\cos\left(\alpha-45^{^\circ}\right)>0$}
{$ C=\tan\left(270^{^\circ}-\alpha\right)<0$}
{$ D=\cos\left(2\alpha+90^{^\circ}\right)>0$}
\loigiai{
\begin{itemchoice}
\itemch  Ta có $0^{\circ}<\alpha<90^{\circ}\Rightarrow 90^{\circ}<\alpha+90^{\circ}<180^{\circ}$ 
$\Rightarrow\sin\left(\alpha+90^{\circ}\right)>0$.
\itemch Ta có $0^{\circ}<\alpha<90^{\circ}\Rightarrow-45^{\circ}<\alpha-45^{\circ}<45^{\circ}$
$\Rightarrow\cos\left(\alpha-45^{\circ}\right)>0$.
\itemch Ta có $0^{\circ}<\alpha<90^{\circ}\Rightarrow-90^{\circ}<-\alpha<0^{\circ}$\\
$\begin{aligned}
&\Rightarrow 270^{\circ}+\left(-90^{\circ}\right)<270^{\circ}+(-\alpha)<270^{\circ}+0^{\circ}\\
&\Rightarrow 180^{\circ}<270^{\circ}-\alpha<270^{\circ}\\
&\Rightarrow\tan\left(270^{\circ}-\alpha\right)>0.
\end{aligned}$
\itemch Ta có $0^{\circ}<\alpha<90^{\circ}\Rightarrow 90^{\circ}<2\alpha+90^{\circ}<270^{\circ}$
$\Rightarrow\cos\left(2\alpha+270^{\circ}\right)<0$.
\end{itemchoice}
}
\end{ex}

\begin{ex}%[1D1H2-2]
Biết $\sin x=-\dfrac{3}{5}$ với $\pi<x<\dfrac{3\pi}{2}$. Các khẳng định  sau đây đúng hay sai? 
\choiceTF
{$\cos x>0$}
{\True $\cos x=-\dfrac{4}{5}$}
{\True $\tan x=\dfrac{3}{4}$}
{\True $\cot x=\dfrac{4}{3}$}
\loigiai{
Do $\pi<x<\dfrac{3\pi}{2}$ nên $\cos x<0$.\\
Ta có $\sin ^2 x+\cos ^2 x=1\Rightarrow\cos ^2 x=1-\sin ^2 x=1-\dfrac{9}{25}=\dfrac{16}{25}$.\\
$\Rightarrow\cos x=-\dfrac{4}{5};\tan x=\dfrac{\sin x}{\cos x}=\dfrac{3}{4};\cot x=\dfrac{\cos x}{\sin x}=\dfrac{4}{3}$.}
\end{ex}



\begin{ex}%[1D1H2-2]
Biết $\tan\alpha=\dfrac{2\sqrt{10}}{9},\pi<\alpha<\dfrac{3\pi}{2}$. Các khẳng định  sau đây đúng hay sai? 
\choiceTF
{\True $\cot\alpha=\dfrac{9}{2\sqrt{10}}$}
{\True $\cos\alpha=-\dfrac{9}{11}$}
{\True $\left\{\begin{array}{*{35}{l}}
\cos\alpha <0\\
\sin\alpha <0\\
\end{array}\right.$}
{$\sin\alpha=-\dfrac{2\sqrt{10}}{11}$}
\loigiai{
$\pi<\alpha<\dfrac{3\pi}{2}\Rightarrow\heva{\cos\alpha<0\\\sin\alpha<0}$; $\cot\alpha=\dfrac{9}{2\sqrt{10}}$; \\
$\cos\alpha=-\sqrt{\dfrac{1}{\tan ^2\alpha+1}}=-\sqrt{\dfrac{1}{\dfrac{40}{81}+1}}=-\dfrac{9}{11}$;\\
$\sin\alpha=-\sqrt{1-\cos ^2\alpha}=-\sqrt{1-\dfrac{81}{121}}=\dfrac{2\sqrt{10}}{11}$.
}
\end{ex}

\Closesolutionfile{ans}
\indapan{3}{ans/ans-0-B1-DS}

\Opensolutionfile{ans}[ans/ans-0-B1-KQ]
\caukq



\begin{ex}%[1D1H2-3]
Biểu thức sau $ T=2\sin\left(\dfrac{9\pi}{2}-x\right)+3\cos (19\pi-x)=k\cos x$. Khi đó $ k$ bằng bao nhiêu?
\shortans{$-1$}
\loigiai{
Ta có \\
$ \begin{aligned}
T&=2\sin\left(4\pi+\dfrac{\pi}{2}-x\right)+3\cos (18\pi+\pi-x)\\
&=2\sin\left(\dfrac{\pi}{2}-x\right)+3\cos (\pi-x)=2\cos x-3\cos x=-\cos x  \Rightarrow k=-1.
\end{aligned}
$
}
\end{ex}



\begin{ex}%[1D1H2-3]
Biểu thức $ A=\tan\left(\dfrac{17\pi}{2}-x\right)+2\cot (5\pi+x)=k\cot x$, Khi đó $ k$ bằng bao nhiêu?
\shortans{$3$}
\loigiai{
$A=\tan\left(8\pi+\dfrac{\pi}{2}-x\right)+2\cot x=\tan\left(\dfrac{\pi}{2}-x\right)+2\cot x=\cot x+2\cot x=3\cot x$.\\
$\Rightarrow k=3.$
}
\end{ex}





\begin{ex}%[1D1H2-3]
Cho $\cos\alpha=\dfrac{3}{4}$. Tính giá trị của biểu thức $B=\dfrac{\tan\alpha+3\cot\alpha}{\tan\alpha+\cot\alpha}$ (làm tròn đến phần trăm).\\
\shortans{$2{,}13$}
\loigiai{
Ta có $\cos\alpha=\dfrac{3}{4}\Rightarrow 1+\tan ^2\alpha=\dfrac{1}{\cos ^2\alpha}=\dfrac{16}{9}\Rightarrow\tan ^2\alpha=\dfrac{7}{9}$.\\
Khi đó $B=\dfrac{\tan\alpha+\dfrac{3}{\tan\alpha}}{\tan\alpha+\dfrac{1}{\tan\alpha}}=\dfrac{\tan ^2\alpha+3}{\tan ^2\alpha+1}=\dfrac{\dfrac{7}{9}+3}{\dfrac{7}{9}+1}=\dfrac{17}{8}=2{,}13$.}
\end{ex}



\begin{ex}%[1D1V2-3]
Biết $\sin a+\cos a=\sqrt{2}$. Tính giá trị của $\sin ^4 a+\cos ^4 a$.\\
\shortans{$0{,}5$}
\loigiai{
Ta có\\
$\begin{aligned}
\sin a+\cos a =\sqrt{2}\Rightarrow 2&=(\sin a+\cos a)^2\\
&=\sin ^2 a+2\sin a\cos a+\cos ^2 a\\
&=1+2\sin a\cos a 
\end{aligned}$\\
$\Rightarrow\sin a\cos a=\dfrac{1}{2}$.\\
Khi đó $\sin ^4 a+\cos ^4 a=\left(\sin ^2 a+\cos ^2 a\right)^2-2\sin ^2 a\cos ^2 a=1-2\left(\dfrac{1}{2}\right)^2=\dfrac{1}{2}=0{,}5$.}
\end{ex}



\begin{ex}%[1D1V2-3]
Cho biểu thức $f(x)=3\left(\sin ^4 x+\cos ^4 x\right)-2\left(\sin ^6 x+\cos ^6 x\right)$. Tính $ f(1)$.
\shortans{$1$}
\loigiai{
Ta có $\sin ^4 x+\cos ^4 x=1-2\sin ^2 x\cos ^2 x$, $\sin ^6 x+\cos ^6 x=1-3\sin ^2 x\cos ^2 x$.\\
Suy ra $f(x)=3\left(1-2\sin ^2 x\cos ^2 x\right)-2\left(1-3\sin ^2 x\cos ^2 x\right)=1$.\\
Vậy biểu thức $f(1) =1$.}
\end{ex}






\begin{ex}%[1D1V2-3]
Cho $\tan x=-\dfrac{4}{3}$ và $\dfrac{\pi}{2}<x<\pi$. Tính giá trị của biểu thức $M=\dfrac{\sin ^2 x-\cos x}{\sin x-\cos ^2 x}$ (làm tròn đến phần trăm).\\
\shortans{$2{,}82$}
\loigiai{
Ta có $\tan x=-\dfrac{4}{3}\Rightarrow\cos ^2 x=\dfrac{1}{1+\tan ^2 x}=\dfrac{9}{25}\Rightarrow\cos x=\pm\dfrac{3}{5}$.\\
Vì $\dfrac{\pi}{2}<x<\pi\Rightarrow\cos x=-\dfrac{3}{5}\Rightarrow\sin x=\tan x\cdot\cos x=\dfrac{4}{5}$ \\
$\Rightarrow M=\dfrac{\sin ^2 x-\cos x}{\sin x-\cos ^2 x}=\dfrac{31}{11}\approx 2{,}82$.}
\end{ex}



\Closesolutionfile{ans}
\indapan{6}{ans/ans-0-B1-KQ}

\subsection{Đề 2}

\Opensolutionfile{ans}[ans/ans-0-B1]
\caulc



\begin{ex}%[1D1H2-4]
	Tìm khẳng định {\bf sai} (với điều kiện các hệ thức đã xác định)?
	\choice
	{$\tan\left(\pi+\alpha\right)=\tan\alpha $}
	{\True $\cos\left(\dfrac{\pi}{2}+\alpha\right)=\sin\alpha $}
	{$\cot\left(-\alpha\right)=-\cot\alpha $}
	{$\sin\left(\pi-\alpha\right)=\sin\alpha $}
	\loigiai{ 
		$\cos\left(\dfrac{\pi}{2}+\alpha\right)=-\sin\alpha $ nên  $\cos\left(\dfrac{\pi}{2}+\alpha\right)=\sin\alpha $ là khẳng định sai.	
	}
\end{ex}



\begin{ex}%[1D1B2-2]%Câu 4
	Cho $\dfrac{\pi}{2}<\alpha <\pi $. Kết quả đúng là
	\choice
	{$\sin\alpha >0,$ $\cos\alpha >0$}
	{$\sin\alpha <0,$ $\cos\alpha <0$}
	{\True $\sin\alpha >0,$ $\cos\alpha <0$}
	{$\sin\alpha <0,$ $\cos\alpha >0$}
	\loigiai
	{	Vì $\dfrac{\pi}{2}<\alpha <\pi $ nên $\alpha $ nằm trong góc phần tư thứ $ 2$ nên $\sin\alpha >0,\cos\alpha <0$.}
\end{ex}



\begin{ex}%[1D1H2-2]
	Cho $\dfrac{7\pi}{4}<\alpha <2\pi $. Khẳng định  nào sau đây đúng?
	\choice
	{\True $\cos\alpha >0$}
	{$\sin\alpha >0$}
	{$\tan\alpha >0$}
	{$\cot\alpha >0$}
	\loigiai
	{$\dfrac{7\pi}{4}<\alpha <2\pi\Leftrightarrow\dfrac{3\pi}{2}+\dfrac{\pi}{4}<\alpha <2\pi $ nên $\alpha$ thuộc cung phần tư thứ $IV$ vì vậy $\cos\alpha >0$.}
\end{ex}



\begin{ex}%[1D1H2-2]
	Cho biết $\tan\alpha=\dfrac{1}{2}$. Tính $\cot\alpha $
	\choice
	{\True $\cot\alpha=2$}
	{$\cot\alpha=\dfrac{1}{4}$}
	{$\cot\alpha=\dfrac{1}{2}$}
	{$\cot\alpha=\sqrt{2}$}
	\loigiai{Ta có $\tan\alpha \cdot \cot\alpha=1 \Rightarrow\cot\alpha=\dfrac{1}{\tan\alpha}=\dfrac{1}{\dfrac{1}{2}}=2$.}
\end{ex}



\begin{ex}%[1D1H2-2]
	Trong các đẳng thức sau, đẳng thức nào đúng?
	\choice
	{$\cos 150^\circ=\dfrac{\sqrt{3}}{2}$}
	{$\cot 150^\circ=\sqrt{3}$}
	{\True $\tan 150^\circ=-\dfrac{1}{\sqrt{3}}$}
	{$\sin 150^\circ=-\dfrac{\sqrt{3}}{2}$}
	\loigiai
	{	$\cos 150^\circ=-\dfrac{\sqrt{3}}{2}$ suy ra $\cos 150^\circ=\dfrac{\sqrt{3}}{2}$ sai.\\
		$\cot 150^\circ=-\sqrt{3}$ suy ra $\cot 150^\circ=\sqrt{3}$ sai.\\
		$\tan 150^\circ=-\dfrac{1}{\sqrt{3}}$ suy ra $\tan 150^\circ=-\dfrac{1}{\sqrt{3}}$ đúng.\\
		$\sin 150^\circ=\dfrac{1}{2}$ suy ra $\sin 150^\circ=-\dfrac{\sqrt{3}}{2}$ sai.}
\end{ex}



\begin{ex}%[1D1H2-2]
	Cho $\sin\alpha=\dfrac{3}{5}\left(\dfrac{\pi}{2}<\alpha <\pi\right)$. Giá trị của $\cos\alpha $ bằng
	\choice
	{\True $\cos\alpha=-\dfrac{4}{5}$}
	{$\cos\alpha=\dfrac{4}{5}$}
	{$\cos\alpha=\dfrac{2}{5}$}
	{$\cos\alpha=-\dfrac{2}{5}$}
	\loigiai{
		Ta có $\sin^2\alpha+\cos^2\alpha=1\Rightarrow{\cos^2}\alpha=1-\sin^2\alpha=1-\dfrac{9}{25}=\dfrac{16}{25}$.\\
		Mà $\dfrac{\pi}{2}<\alpha <\pi\Rightarrow\cos\alpha=-\dfrac{4}{5}$.}
\end{ex}



\begin{ex}%[1D1H2-3]
	Đơn giản biểu thức $ A=\cos\left(\alpha-\dfrac{\pi}{2}\right)$, ta được
	\choice
	{$\cos\alpha $}
	{\True $\sin\alpha $}
	{$\cos\alpha $}
	{$-\sin\alpha $}
	\loigiai{Ta có $ A=\cos\left(\alpha-\dfrac{\pi}{2}\right)$$=\cos\left(\dfrac{\pi}{2}-\alpha\right)$ $=\sin\alpha $.}
\end{ex}



\begin{ex}%[1D1H2-3]
	Cho $ a$ là số thực bất kỳ. Chọn khẳng định đúng trong các khẳng định sau:
	\choice
	{$\sin a+\cos a=1$}
	{$\sin^3a+\cos^3a=1$}
	{$\sin^4a+\cos^4a=1$}
	{\True $\sin^2a+\cos^2a=1$}
	\loigiai{Ta có $\sin^2a+\cos^2a=1$ với mọi $a\in\mathbb{R}$ nên $\sin^2a+\cos^2a=1$ đúng.}
\end{ex}




\begin{ex}%[1D1H2-3]
	Tính $S=\sin^25^0+\sin^210^0+\sin^215^0+ \cdots \sin^28^0+\sin^285^0$ 
	\choice
	{$\dfrac{19}{2}$}
	{$ 8$}
	{\True $\dfrac{17}{2}$}
	{$ 9$}
	\loigiai
	{	$\begin{aligned}
		S&=\sin^25^0+\sin^210^0+\sin^215^0+\cdots +\sin^280^0+\sin^285^0\\
		&=\left(\sin^25^0+\sin^285^0\right)+\left(\sin^210^0+\sin^280^0\right)+ \cdots + \left(\sin^215^0+\sin^275^0\right)+\sin^245^0\\
		&=8+\dfrac{1}{2}=\dfrac{17}{2}.
		\end{aligned}
	$
		}
\end{ex}





\begin{ex}%[1D1H2-3]
	Biết $\cos\alpha=-\dfrac{1}{3}$. Khi đó, giá trị biểu thức $ A=\dfrac{5\cot\alpha+4\tan\alpha}{5\cot\alpha-4\tan\alpha}$ là
	\choice
	{\True $-\dfrac{37}{27}$}
	{$\dfrac{37}{27}$}
	{$\dfrac{15}{27}$}
	{$\dfrac{19}{27}$}
	\loigiai
	{Ta có $ 1+\tan^2\alpha=\dfrac{1}{\cos^2\alpha}\Rightarrow{\tan^2}\alpha=\dfrac{1}{\cos^2\alpha}-1=\dfrac{1}{\left(-\dfrac{1}{3}\right)^2}-1=8$.\\
		$ A=\dfrac{5\cot\alpha+4\tan\alpha}{5\cot\alpha-4\tan\alpha}=\dfrac{5+4\tan^2\alpha}{5-4\tan^2\alpha}=\dfrac{5+4.8}{5-4.8}=-\dfrac{37}{27}$.}
\end{ex}





\begin{ex}%[1D1V2-3]
	Tính giá trị biểu thức $P=\sin^2\dfrac{\pi}{6}+\sin^2\dfrac{\pi}{3}+\sin^2\dfrac{\pi}{4}+\sin^2\dfrac{9\pi}{4}+\tan\dfrac{\pi}{6}\cot\dfrac{\pi}{6}$ .
	\choice
	{$2$}
	{$4$}
	{\True $3$}
	{$1$}
	\loigiai
	{Ta có \\
		$\sin\dfrac{\pi}{6}=\dfrac{1}{2}\Rightarrow{\sin^2}\dfrac{\pi}{6}=\dfrac{1}{4}$, \\ $\sin\dfrac{\pi}{3}=\dfrac{\sqrt{3}}{2}\Rightarrow{\sin^2}\dfrac{\pi}{3}=\dfrac{3}{4}$, \\ $\sin\dfrac{\pi}{4}=\dfrac{\sqrt{2}}{2}\Rightarrow{\sin^2}\dfrac{\pi}{4}=\dfrac{1}{2}$,\\
		$\sin\dfrac{9\pi}{4}=\sin\left(\dfrac{\pi}{4}+2\pi\right)=\sin\dfrac{\pi}{4}=\dfrac{\sqrt{2}}{2}$ $\Rightarrow{\sin^2}\dfrac{9\pi}{4}=\dfrac{1}{2}$, $\tan\dfrac{\pi}{6}\cdot \cot\dfrac{\pi}{6}=1$.\\
		Suy ra $P=\dfrac{1}{4}+\dfrac{3}{4}+\dfrac{1}{2}+\dfrac{1}{2}+1=3$.}
\end{ex}



\begin{ex}%[1D1K2-3]
	Giá trị lớn nhất của biểu thức $ T=\cos^2x-2\sin^2x$là
	\choice
	{$-2$}
	{$ 4$}
	{\True $ 1$}
	{$ 2$}
	\loigiai
	{		$ T=1-\sin^2x-2\sin^2x$$=1-3\sin^2x$.\\
		Ta có $ 0\le{\sin^2}x\le 1$$\Leftrightarrow 0\ge-3\sin^2x\ge-3,\forall x\in R$$\Leftrightarrow 1\ge 1-3\sin^2x\ge-2$.\\
		Giá trị lớn nhất là $ 1$.}
\end{ex}

\Closesolutionfile{ans}
\indapan{6}{ans/ans-0-B1}
\cauds
\Opensolutionfile{ans}[ans/ans-0-B1-DS]

\begin{ex}%[1D1V2-3]
	Các khẳng định  sau đây đúng hay sai? 
	\choiceTF
	{\True $ A=\dfrac{1}{\sqrt{3}}\cos{30^{^\circ}}-3\sqrt{2}\sin{45^{^\circ}}+\cot{45^{^\circ}}=-\dfrac{3}{2}$}
	{\True $ B=\sin^260^{^\circ}+\tan^230^{^\circ}-2=-\dfrac{11}{12}$}
	{\True $ C=2\sin\dfrac{\pi}{6}+\sqrt{3}\cos\dfrac{\pi}{6}=\dfrac{5}{2}$}
	{$ D=\dfrac{\tan\dfrac{\pi}{4}+1}{\dfrac{1}{\sqrt{2}}\cot\dfrac{\pi}{2}-2}=\dfrac{6}{5}$}
	\loigiai{
		\begin{itemchoice}		
			\itemch $A=\dfrac{1}{\sqrt{3}}\cdot\dfrac{\sqrt{3}}{2}-3\sqrt{2}\cdot\dfrac{\sqrt{2}}{2}+1=\dfrac{1}{2}-3+1=-\dfrac{3}{2}$.
			\itemch $B=\left(\dfrac{\sqrt{3}}{2}\right)^2+\left(\dfrac{1}{\sqrt{3}}\right)^2-2=\dfrac{3}{4}+\dfrac{1}{3}-2=-\dfrac{11}{12}$.
			\itemch $C=2\cdot\dfrac{1}{2}+\sqrt{3}\cdot\dfrac{\sqrt{3}}{2}=\dfrac{5}{2}$.
			\itemch $D=\dfrac{1+1}{\dfrac{1}{\sqrt{3}}\cdot\dfrac{1}{\sqrt{3}}-2}=-\dfrac{6}{5}$.
		\end{itemchoice}
	}
\end{ex}


\begin{ex}%[1D1H2-2]
	Cho $\alpha=-\dfrac{\pi}{4}+(2k+1)\pi $ (biết $k\in\mathbb{Z}$). Các khẳng định  sau đây đúng hay sai? 
	\choiceTF
	{\True $\sin\alpha=\dfrac{\sqrt{2}}{2}$}
	{\True $\cos\alpha=\dfrac{\sqrt{2}}{2}$}
	{$\tan\alpha=-1$}
	{$\cot\alpha=-1$}
	\loigiai{
		Ta có \\
		$\begin{aligned}
			\sin\left[-\dfrac{\pi}{4}+(2 k+1)\pi\right]& =\sin\left(-\dfrac{\pi}{4}+2 k\pi+\pi\right)=\sin\left(-\dfrac{\pi}{4}+\pi\right) \\
			&=-\sin\left(-\dfrac{\pi}{4}\right)=\sin\dfrac{\pi}{4}=\dfrac{\sqrt{2}}{2}.
			\end{aligned}$\\
			$\begin{aligned}
			\cos\left[-\dfrac{\pi}{4}+(2k+1)\pi\right]& =\cos\left(-\dfrac{\pi}{4}+2k\pi+\pi\right)=\cos\left(-\dfrac{\pi}{4}+\pi\right)\\ 
			&=-\cos\left(-\dfrac{\pi}{4}\right)=-\cos\dfrac{\pi}{4}=-\dfrac{\sqrt{2}}{2}.
		\end{aligned}$\\
				$\begin{aligned}
			&\tan\left[-\dfrac{\pi}{4}+(2k+1)\pi\right]=\tan\left(-\dfrac{\pi}{4}\right)=-\tan\dfrac{\pi}{4}=-1.\\ 
			&\cot\left[-\dfrac{\pi}{4}+(2k+1)\pi\right]=\cot\left(-\dfrac{\pi}{4}\right)=-\cot\dfrac{\pi}{4}=-1.
		\end{aligned}$
	}
\end{ex}

\begin{ex}%[1D1H2-3]
	Cho biết $\sin\alpha=\dfrac{3}{5},\cos\alpha=-\dfrac{4}{5}$ và các biểu thức \\ 
	\centerline{$ A=\sin\left(\dfrac{\pi}{2}-\alpha\right)+\sin (\pi+\alpha)$; $ B=\cos (\pi-\alpha)+\cot\left(\dfrac{\pi}{2}-\alpha\right)$.} \\
	Các khẳng định  sau đây đúng hay sai? 
	\choiceTF
	{\True $ A=\cos\alpha-\sin\alpha $}
	{$ B=\cos\alpha+\tan\alpha $}
	{$ A+B=\dfrac{27}{20}$}
	{\True $ A-B=-\dfrac{29}{20}$}
	\loigiai{
		Ta có $ A=\sin\left(\dfrac{\pi}{2}-\alpha\right)+\sin (\pi+\alpha)=\cos\alpha-\sin\alpha=-\dfrac{4}{5}-\dfrac{3}{5}=-\dfrac{7}{5}$.\\
		Ta có $ B=\cos (\pi-\alpha)+\cot\left(\dfrac{\pi}{2}-\alpha\right)=-\cos\alpha+\tan\alpha $.
		$=-\cos\alpha+\dfrac{\sin\alpha}{\cos\alpha}=\dfrac{4}{5}+\dfrac{\dfrac{3}{5}}{-\dfrac{4}{5}}=\dfrac{1}{20}.$}
\end{ex}


\begin{ex}%[1D1V2-5]
	Từ một vị trí ban đầu trong không gian, vệ tinh $X$ chuyển động theo quỹ đạo là một đường tròn quanh Trái Đất và luôn cách tâm Trái Đất một khoảng bằng $9200$ km. Sau 2 giờ thì vệ tinh $X$ hoàn thành hết một vòng di chuyển. Các khẳng định  sau đây đúng hay sai? 
	\choiceTF
	{\True Quãng đường vệ tinh $X$ chuyển động được sau $1$ giờ là $\approx 28902{,}65$ (km)}
	{\True Quãng đường vệ tinh $X$ chuyển động được sau $1{,}5$ giờ là $\approx 43353{,}98$ (km)}
	{Sau khoảng 5,3 giờ thì $X$ di chuyển được quãng đường $240000$ (km)}
	{\True Giả sử vệ tinh di chuyển theo chiều dương của đường tròn, sau $4{,}5$ giờ thì vệ tinh vẽ nên một góc $\dfrac{9\pi}{2}$ rad}
	\loigiai{
		\begin{itemchoice}
			\itemch Một vòng di chuyển của $X$ chính là chu vi đường tròn\\
			\centerline{$C=2\pi R=2\pi .9200=18400\pi$ (km).}\\
			Sau 1 giờ, vệ tinh di chuyển nửa đường tròn với quãng đường là \\
			\centerline{$\dfrac{1}{2}C=9200\pi\approx 28902,65$ (km).}
			\itemch Sau 1,5 giờ, vệ tinh di chuyển được $\dfrac{1,5.1}{2}$ đường tròn (hay $\dfrac{3}{4}$ đường tròn), quãng đường là\\ \centerline{$\dfrac{3}{4}C=\dfrac{3}{4}\cdot 18400\pi=13800\pi\approx 43353,98$ (km).}
			\itemch Số giờ để vệ tinh $X$ thực hiện quãng đường $240000$ km là $\dfrac{240000}{9200\pi}\approx 8,3$ (giờ).
			\itemch Sau $4{,}5$ giờ thì số vòng tròn mà vệ tinh $X$ di chuyển được là $\dfrac{4{,}5}{2}=\dfrac{9}{4}$ (vòng).
			\itemch Số đo góc lượng giác thu được là $\dfrac{9}{4}\cdot 2\pi=\dfrac{9\pi}{2}$ (rad).
		\end{itemchoice}
	}
\end{ex}




	\Closesolutionfile{ans}
	\indapan{3}{ans/ans-0-B1-DS}
	
	\Opensolutionfile{ans}[ans/ans-0-B1-KQ]
	\caukq
	
\begin{ex}%[1D1H2-3]
	Cho $\cos x=\dfrac{1}{2}$ . Tính giá trị biểu thức $P=3\sin ^2 x+4\cos ^2 x$.
	\shortans{$3{,}25$}	
	\loigiai{
		Ta có $\cos x=\dfrac{1}{2}\Rightarrow\sin ^2 x=1-\cos ^2 x=1-\dfrac{1}{4}=\dfrac{3}{4}$.\\
		Khi đó $P=3\sin ^2 x+4\cos ^2 x=3\cdot\dfrac{3}{4}+4\cdot\dfrac{1}{4}=\dfrac{13}{4}=3{,}25$.}
\end{ex}





\begin{ex}%[1D1V2-3]
	Cho tam giác $A B C$, khi đó biểu thức\\ \centerline{$\dfrac{\sin^3\dfrac{B}{2}}{\cos\left(\dfrac{A+C}{2}\right)}+\dfrac{\cos^3\dfrac{B}{2}}{\sin\left(\dfrac{A+C}{2}\right)}-\dfrac{\cos (A+C)}{\sin B}\cdot \tan B$}\\ bằng?
	\shortans{2}
	\loigiai{
		Vì $A+B+C=180^{\circ}$ nên $A+C=180^{\circ}-B$.\\
		$\begin{aligned}
			\text{VT}& =\dfrac{\sin^3\dfrac{B}{2}}{\cos\left(\dfrac{180^{^\circ}-B}{2}\right)}+\dfrac{\cos^3\dfrac{B}{2}}{\sin\left(\dfrac{180^{^\circ}-B}{2}\right)}-\dfrac{\cos\left(180^\circ-B\right)}{\sin B}\cdot\tan B\\ 
			&=\dfrac{\sin^3\dfrac{B}{2}}{\sin\dfrac{B}{2}}+\dfrac{\cos^3\dfrac{B}{2}}{\cos\dfrac{B}{2}}-\dfrac{-\cos B}{\sin B}\cdot\tan B\\ 
			&=\sin^2\dfrac{B}{2}+\cos^2\dfrac{B}{2}+\cot B\cdot\tan B=1+1=2=\text{VP}.
		\end{aligned}$}
\end{ex}




\begin{ex}%[1D1H2-3]
	Cho $\cot\alpha=\dfrac{1}{3}$. Tính giá trị của biểu thức $A=\dfrac{3\sin\alpha+4\cos\alpha}{2\sin\alpha-5\cos\alpha}$.
	\shortans{13}
	\loigiai{
		Vì $\cot\alpha=\dfrac{1}{3}$ nên $\sin\alpha\neq 0$. Chia cả tử và mẫu của biểu thức $A$ cho $\sin\alpha$, ta có\\ $A=\dfrac{3+4\cot\alpha}{2-5\cot\alpha}=\dfrac{3+4\cdot\dfrac{1}{3}}{2-5\cdot\dfrac{1}{3}}=13$.}
\end{ex}



\begin{ex}%[1D1H2-3]
	Cho $\tan\alpha=\sqrt{2}$. Tính giá trị của biểu thức $C=\dfrac{\sin\alpha-\cos\alpha}{\sin ^3\alpha+3\cos ^3\alpha+2\sin\alpha}$ (làm tròn đến phần trăm).
	\shortans{$0{,}09$}
	\loigiai{
		Vì $\tan\alpha=\sqrt{2}$ nên $\cos\alpha\neq 0$. Chia cả tử và mẫu của biểu thức $C$ cho $\cos ^3\alpha$, ta được\\
		$\begin{aligned}
			 C&=\dfrac{\dfrac{\sin\alpha}{\cos\alpha}\cdot\dfrac{1}{\cos^2\alpha}-\dfrac{1}{\cos^2\alpha}}{\tan^3\alpha+3+2\cdot\dfrac{\sin\alpha}{\cos\alpha}\cdot\dfrac{1}{\cos^2\alpha}}=\dfrac{\tan\alpha\left(1+\tan^2\alpha\right)-\left(1+\tan^2\alpha\right)}{\tan^3\alpha+3+2\tan\alpha\left(1+\tan^2\alpha\right)}\\ 
			&=\dfrac{\left(1+\tan^2\alpha\right)(\tan\alpha-1)}{3\tan^3\alpha+2\tan\alpha+3}=\dfrac{(1+2)(\sqrt{2}-1)}{3\cdot 2\sqrt{2}+2\sqrt{2}+3}=\dfrac{3(\sqrt{2}-1)}{8\sqrt{2}+3} \approx 0{,}09.
		\end{aligned}$}
\end{ex}



\begin{ex}%[1D1K2-3]
	Cho $\tan\alpha+\cot\alpha=m$. Tìm $m>0$ để $\tan ^2\alpha+\cot ^2\alpha=7$.\\
	\shortans{$ 3$}
	\loigiai{
		Ta có $7=\tan ^2\alpha+\cot ^2\alpha=(\tan\alpha+\cot\alpha)^2-2\tan\alpha\cot\alpha=m^2-2$.\\
		Suy ra $m^2=9\Rightarrow m=3$.}
\end{ex}








\begin{ex}%[1D1H2-3]
	Cho hai góc nhọn $a$ và $b$. Biết $\cos a=\dfrac{1}{3}$ và $\cos b=\dfrac{1}{4}$. Tính giá trị của
	$P=(\cos a\cdot\cos b)^2-(\sin a\cdot\sin b)^2 .$ (làm tròn đến chục)
	\shortans{$-0{,}8$}
	\loigiai{
		Ta có\\
		$\begin{aligned}
			{P}&=(\cos a\cos b)^2-(\sin a\sin b)^2=(\cos a\cos b)^2-\left(1-\cos ^2 a\right)\left(1-\cos ^2 b\right)\\
			&=\left(\dfrac{1}{12}\right)^2-\dfrac{8}{9}\cdot\dfrac{15}{16}=-\dfrac{119}{144}\approx-0{,}8.
		\end{aligned}$}
\end{ex}


	
	\Closesolutionfile{ans}
	\indapan{6}{ans/ans-0-B1-KQ}